% !TEX root = ../main.tex
\section{Vorbereitung (Aufgabe 5)}
\label{sec:vorbereitung}

\subsection{Stromverlauf bei Grundfrequenztaktung}

Verlangt ist die Skizze des Stromverlaufs bei Grundfrequenztaktung unter der Annahme, dass die Grundschwingung der Wechselrichterspannung mit der Netzspannung übereinstimmt ($U_{\mathrm{WR1}}=\UN$, $\Th=0$). Die Spannung über der Induktivität ist dann $\uWR-\uN$, und es gilt $L\,\mathrm{d}i/\mathrm{d}t=\uWR-\uN$. Für die erste Halbperiode $0\le\omega t\le\pi$ mit $\uWR=+\Udc$ und $\uN=\hat U_{\mathrm N}\sin\omega t$ folgt durch Integration
\begin{equation}
  i(\omega t)=\frac{1}{X}\Bigl[\Udc\bigl(\omega t-\tfrac{\pi}{2}\bigr)+\hat U_{\mathrm N}\cos\omega t\Bigr],
  \qquad i(\omega t+\pi)=-i(\omega t).
  \label{eq:strom_vorb}
\end{equation}
Die Integrationskonstante ist so gewählt, dass der Strom mittelwertfrei ist. Mit $\hat U_{\mathrm N}=\sqrt2\cdot\SI{230}{\volt}=\SI{325.3}{\volt}$, $\Udc=\tfrac{\pi}{4}\hat U_{\mathrm N}=\SI{255.5}{\volt}$ und $X=\SI{22.0}{\ohm}$ ergibt sich der Verlauf in Abbildung~\ref{fig:strom_vorb} mit einem Scheitelwert von etwa \SI{3.5}{\ampere}. Der Strom besteht nur aus Oberschwingungen; seine Form entspricht dem Integral von «Rechteck minus Sinus».

\begin{figure}[htbp]
  \centering
  \begin{tikzpicture}
    \begin{axis}[rep,height=4.6cm,xlabel={$\omega t$ in \si{\degree}},ylabel={$u$ in V},
        xmin=0,xmax=360,xtick={0,90,...,360},ymin=-400,ymax=400,
        legend style={at={(0.5,1.04)},anchor=south,legend columns=2,font=\small}]
      \addplot[thick,domain=0:360,samples=721]{ifthenelse(x<180,255.46,-255.46)};
      \addlegendentry{$\uWR$ (Rechteck)}
      \addplot[thick,dashed,domain=0:360,samples=361]{325.27*sin(x)};
      \addlegendentry{$\uN$}
    \end{axis}
  \end{tikzpicture}

  \begin{tikzpicture}
    \begin{axis}[rep,height=4.6cm,xlabel={$\omega t$ in \si{\degree}},ylabel={$i$ in A},
        xmin=0,xmax=360,xtick={0,90,...,360},ymin=-5,ymax=5]
      \addplot[thick,domain=0:360,samples=721]{ifthenelse(x<180,(255.46*(x*pi/180-pi/2)+325.27*cos(x))/21.99,-(255.46*((x-180)*pi/180-pi/2)+325.27*cos(x-180))/21.99)};
    \end{axis}
  \end{tikzpicture}
  \caption{Berechnete Spannungen (oben) und Strom (unten) bei Grundfrequenztaktung mit $U_{\mathrm{WR1}}=\UN$ und $\Th=0$, nach \eqref{eq:strom_vorb}.}
  \label{fig:strom_vorb}
\end{figure}

\subsection{Beeinflussung von Wirk- und Blindleistung}

Abbildung~\ref{fig:zeiger} zeigt das Zeigerdiagramm mit dem Spannungsabfall $\Delta\underline U=\underline U_{\mathrm N}-\underline U_{\mathrm{WR1}}=\mathrm{j}X\underline I$ an der Kopplungsinduktivität; der Strom eilt $\Delta\underline U$ um \SI{90}{\degree} nach. Die Leistungen folgen aus \eqref{eq:PQ}:

\begin{itemize}
  \item Die \emph{Wirkleistung} wird über den Winkel $\Th$ zwischen $\underline U_{\mathrm{WR1}}$ und $\underline U_{\mathrm N}$ eingestellt, $P\propto\sin\Th$. Am Bedienfeld geschieht dies mit den Tasten für $\Th$.
  \item Die \emph{Blindleistung} wird über die Amplitude $U_{\mathrm{WR1}}$ eingestellt. Bei fester Pulsform ist $U_{\mathrm{WR1}}\propto\Udc$, am Versuch also über den Erregerstrom der GM. Ist $U_{\mathrm{WR1}}\cos\Th<\UN$, bezieht der Wechselrichter induktive Blindleistung aus dem Netz (untererregt); ist $U_{\mathrm{WR1}}\cos\Th>\UN$, liefert er induktive Blindleistung ans Netz (übererregt).
\end{itemize}

\begin{figure}[htbp]
  \centering
  \begin{tikzpicture}[>={Stealth[length=2.5mm]},scale=1.15]
    \coordinate (O) at (0,0);
    \coordinate (UN) at (4,0);
    \coordinate (UW) at (-25:3.6);
    \draw[->,thick] (O) -- (UN) node[right] {$\underline U_{\mathrm N}$};
    \draw[->,thick] (O) -- (UW) node[below right] {$\underline U_{\mathrm{WR1}}$};
    \draw[->,thick,dashed] (UW) -- (UN) node[midway,right,xshift=2pt] {$\Delta\underline U=\mathrm{j}X\underline I$};
    \draw[->,thick] (O) -- (1.5,-0.72) node[below right,xshift=-2pt] {$\underline I$};
    \draw (1.0,0) arc (0:-25:1.0);
    \node at (1.35,-0.2) {$\Th$};
  \end{tikzpicture}
  \caption{Zeigerdiagramm des einphasigen Wechselrichters am Netz bei $P>0$ (Bezug aus dem Netz); die Länge von $\underline I$ ist nicht massstäblich.}
  \label{fig:zeiger}
\end{figure}
